% =============================================================================
% RESEARCH TEAM OVERVIEW
% =============================================================================
%
% COMPILER: pdfLaTeX
%
% FORMATTING:
%   - Single-spaced, 12pt font
%   - 1-inch margins on all sides
%   - Unlocked, searchable PDF
%
% CUSTOM COMMANDS:
%   \sect{Title}    - Section header (garnet bold small caps with period)
%   \subsect{Title} - Subsection header (garnet italic with period)
%
% =============================================================================

\documentclass[12pt]{article}

% --- Page Layout (1-inch margins) ---
\usepackage[margin=1in]{geometry}

% --- Font Encoding & Bold Small Caps ---
\usepackage[T1]{fontenc}
\usepackage{bold-extra}

% --- Colors (UofSC Brand) ---
\usepackage{xcolor}
\definecolor{garnet}{HTML}{73000A}
\definecolor{congaree}{HTML}{1F414D}
\definecolor{atlantic}{HTML}{466A9F}

% --- Paragraph Formatting ---
\setlength{\parindent}{0.25in}
\setlength{\parskip}{0.5em}

% --- Custom Section Commands ---
\newcommand{\sect}[1]{\noindent{\color{garnet}\textbf{\textsc{#1}}.}}
\newcommand{\subsect}[1]{\noindent{\color{garnet}\textit{#1.}}}

% =============================================================================
\begin{document}

% --- Title Block ---
\noindent{\color{garnet}\large\textbf{\textsc{Research Team Overview: Autonomous Systems \& Multimodal Sensor Fusion Laboratory}}}\\[-0.8em]
{\color{garnet}\rule{\textwidth}{0.5pt}}\\[0.2em]
\sect{General Focus} Our research team focuses on sensor validation and multimodal sensor fusion to develop fully autonomous systems capable of operating in novel, unseen environments. We employ a combination of machine learning, neural networks, and reinforcement learning techniques to enable robust perception, decision-making, and control across diverse operational domains.

The team operates primarily in maritime and ground-based testing environments, developing technologies that serve as analogs for GPS-denied, communication-limited, and unstructured scenarios. Our work spans the full autonomy stack---from raw sensor data acquisition through high-level mission planning---with applications in defense, environmental monitoring, and exploration.

Our research methodology emphasizes real-world deployment and validation. We design, build, and test complete autonomous systems rather than isolated algorithmic components. This systems-level approach ensures that our solutions address practical challenges including sensor noise, computational constraints, environmental variability, and integration complexity.

% --- Research Structure ---
\sect{Research Structure} The team is organized into seven primary sub-areas, each containing multiple specialized focus areas. Several topics---including edge deployment, synthetic data generation, and sensor fusion---span multiple sub-areas and represent cross-cutting concerns addressed collaboratively.

\noindent\textbf{Table 1: Research Sub-Areas and Focus Topics}\\[0.2em]
\noindent\begin{tabular}{@{}p{0.28\textwidth}p{0.67\textwidth}@{}}
\hline
\textbf{Sub-Area} & \textbf{Sub-Sub Areas} \\
\hline
Hardware \& Sensors & Sensor selection, physical setups, characterization, calibration, sensor placement, sensor fusion (physical)*, synthetic data generation*, edge deployment* \\[0.3em]
Data Pipeline & Processing, labeling (manual/automated), dataset management, synthetic data generation*, edge deployment* \\[0.3em]
Perception \& State Est. & ML perception, localization, SLAM, sensor fusion (algorithmic)*, edge deployment* \\[0.3em]
Autonomy \& Planning & Path planning, navigation, RL-based control, mission-level decisions \\[0.3em]
Generalization \& Adapt. & Domain adaptation, transfer learning, robustness to novel environments \\[0.3em]
Simulation \& V\&V & Simulation environments, testing, validation, failure analysis \\[0.3em]
Integration & System integration, communications/networking, full stack assembly \\
\hline
\end{tabular}\\[0.3em]
{\small *Cross-cutting topics that span multiple sub-areas}

% =============================================================================
% HARDWARE & SENSORS
% =============================================================================
\sect{Hardware \& Sensors} The Hardware \& Sensors sub-area establishes the physical foundation upon which all downstream processing depends. This team is responsible for selecting, configuring, characterizing, and calibrating the sensors that capture raw environmental data. The quality and reliability of autonomous system performance is fundamentally constrained by the fidelity of sensor measurements, making this sub-area critical to overall system success.

\subsect{Sensor Selection} Sensor selection involves evaluating and choosing appropriate sensing modalities for specific mission requirements. This process considers the operational environment, required measurement accuracy, and practical constraints. The team evaluates sensors across multiple modalities including LiDAR systems for precise range measurements, cameras for rich visual information, radar for weather-robust detection, sonar for underwater applications, inertial measurement units (IMUs) for motion estimation, and GPS receivers for global positioning when available. Selection decisions must balance competing factors including size, weight, and power consumption (SWaP constraints), cost and commercial availability, data output rates and interfaces, environmental robustness, and measurement accuracy under expected operating conditions.

\subsect{Physical Setups} Physical setup design encompasses the mechanical and structural aspects of sensor integration onto autonomous platforms. This includes designing mounting brackets and enclosures that provide rigid attachment while isolating sensors from platform vibrations. The team develops multi-sensor rigs that position complementary sensors for optimal coverage and overlap. Weatherproofing and ruggedization are essential for field deployment, requiring careful attention to sealing, drainage, and thermal management. Platform integration varies significantly across ground vehicles, maritime vessels, and aerial systems, each presenting unique mounting challenges and constraints.

\subsect{Sensor Modeling \& Characterization} Understanding sensor behavior is essential for developing algorithms that appropriately weight and combine measurements. The characterization team develops statistical models of sensor noise, identifying both systematic biases and random error distributions. This work establishes the operating envelope for each sensor, documenting performance degradation under challenging conditions such as low light, adverse weather, or extreme temperatures. Degradation and failure mode analysis identifies how sensors behave as they approach failure, enabling predictive maintenance and graceful degradation strategies.

\subsect{Calibration} Calibration ensures that sensor measurements accurately reflect physical reality. Intrinsic calibration determines per-sensor parameters such as camera focal lengths and lens distortion coefficients or LiDAR beam angles and timing offsets. The team develops and maintains calibration targets including checkerboard patterns for cameras, corner reflectors for radar, and specialized targets for other sensor modalities. Extrinsic calibration establishes the spatial relationships between multiple sensors, enabling measurement fusion in a common coordinate frame. Temporal calibration addresses synchronization between sensors with different sampling rates and latencies.

\subsect{Physical Sensor Placement} Optimizing sensor positions requires balancing multiple objectives including field-of-view coverage, measurement overlap for redundancy, interference avoidance between sensors, and practical accessibility for maintenance. The team uses simulation tools to evaluate candidate configurations before physical installation. Placement decisions must account for platform-specific constraints including available mounting locations, cable routing paths, and interference from platform structures.

\subsect{Sensor Fusion (Physical)} The physical aspects of sensor fusion involve hardware-level considerations for combining data from multiple sensors. This includes designing synchronization systems that trigger sensors simultaneously or with known timing offsets, enabling temporal alignment of measurements. Hardware timestamping with common clock references provides the precision necessary for accurate fusion. Establishing shared coordinate frames requires careful measurement of sensor positions and orientations.

\subsect{Synthetic Data Generation (Hardware)} Hardware contributions to synthetic data include developing accurate sensor models for simulation environments. This requires characterizing not just ideal sensor behavior but also realistic noise, artifacts, and failure modes. Physics-based rendering techniques simulate sensor responses to virtual environments with sufficient fidelity for algorithm training and validation. Hardware-in-the-loop testing connects physical sensors to simulated environments, validating that real hardware performs as expected.

\subsect{Edge/Embedded Deployment (Hardware)} Hardware aspects of edge deployment include selecting embedded compute platforms that provide sufficient processing capability within power and thermal constraints. The team evaluates options including embedded GPUs, FPGAs, and specialized AI accelerators for different application requirements. Power management and thermal design ensure reliable operation in field conditions where active cooling may be limited.

% =============================================================================
% DATA PIPELINE
% =============================================================================
\sect{Data Pipeline} The Data Pipeline sub-area transforms raw sensor outputs into organized, annotated datasets suitable for machine learning training and algorithm validation. This team manages the flow of data from initial capture through processing, labeling, and storage, maintaining the infrastructure that enables all learning-based research activities.

\subsect{Processing} Data processing applies signal conditioning and format standardization to raw sensor outputs. This includes noise filtering to remove sensor artifacts, format conversion to standard representations, and temporal alignment to establish correspondence between asynchronous sensor streams. Compression and storage optimization balance data fidelity against practical storage and transmission constraints. The team develops processing pipelines that can operate both offline for batch processing and online for real-time applications.

\subsect{Labeling (Manual \& Automated)} Creating ground truth annotations for supervised learning is often the most labor-intensive aspect of machine learning development. The labeling team develops workflows and tools that maximize annotator efficiency while maintaining quality standards. This includes designing annotation interfaces tailored to specific tasks such as bounding box annotation, semantic segmentation, or trajectory labeling. Semi-automated labeling pipelines use model predictions to generate initial annotations that human annotators refine, significantly reducing labeling time. Active learning techniques identify the most informative samples for human annotation, focusing labeling effort where it provides greatest value.

\subsect{Dataset Management} Managing large-scale datasets requires robust infrastructure for storage, versioning, and access. The team implements data versioning systems that track dataset evolution over time, enabling reproducible experiments and rollback when issues are discovered. Storage infrastructure balances access speed, capacity, and cost across hot, warm, and cold storage tiers. Dataset splitting into training, validation, and test partitions requires careful attention to avoid data leakage and ensure representative evaluation.

\subsect{Synthetic Data Generation (Data)} Data-side contributions to synthetic data generation include developing augmentation pipelines that expand dataset diversity through geometric transformations, photometric variations, and noise injection. Domain randomization techniques vary simulation parameters to encourage models that generalize beyond specific rendering choices. Procedural scene generation creates diverse virtual environments for training data collection at scale. The team analyzes sim-to-real gaps to identify aspects of synthetic data that fail to transfer to real-world performance.

\subsect{Edge/Embedded Deployment (Data)} Data pipeline considerations for embedded systems include optimizing preprocessing operations for limited computational resources. This may involve reduced-precision arithmetic, algorithmic simplifications, or preprocessing stages that run on specialized hardware. Bandwidth-constrained data transmission requires intelligent filtering and compression to communicate relevant information within communication limits.

% =============================================================================
% PERCEPTION & STATE ESTIMATION
% =============================================================================
\sect{Perception \& State Estimation} The Perception \& State Estimation sub-area transforms processed sensor data into semantic understanding of the environment and estimates of the system's own state. This team develops the machine learning models and estimation algorithms that enable autonomous systems to understand where they are and what surrounds them.

\subsect{ML Perception} Machine learning perception encompasses the neural network models that extract semantic information from sensor data. Object detection and classification identify and categorize entities in the environment such as vehicles, pedestrians, obstacles, and landmarks. Semantic and instance segmentation provide pixel-level or point-level labeling of sensor data. Depth estimation from monocular or stereo imagery provides range information complementing or substituting for active range sensors. Scene understanding synthesizes individual detections into coherent environmental models, reasoning about object relationships, occluded regions, and scene context.

\subsect{Localization} Localization determines the autonomous system's position and orientation within its operating environment. GPS-based localization provides global position when satellite signals are available, though accuracy and availability vary with environment. Visual odometry estimates motion from camera imagery, integrating incremental estimates to track position over time. Inertial navigation uses IMU measurements to estimate motion at high rates, though drift accumulation limits standalone accuracy. Map-based localization matches current observations against prior maps to determine position.

\subsect{SLAM} Simultaneous Localization and Mapping (SLAM) jointly estimates system position and environmental structure when prior maps are unavailable. Visual SLAM uses camera imagery to track features and build sparse or dense environmental models. LiDAR SLAM processes point cloud data for precise geometric mapping. Multi-sensor SLAM fuses information from multiple sensor modalities, exploiting complementary strengths. Loop closure detection and map optimization correct accumulated drift when the system revisits previously observed locations.

\subsect{Sensor Fusion (Algorithmic)} Algorithmic sensor fusion combines information from multiple sensors to achieve estimates superior to any single sensor alone. Early fusion architectures combine raw sensor data before feature extraction, while late fusion combines independent per-sensor predictions. Classical approaches including Kalman filtering and Bayesian estimation provide principled frameworks for combining uncertain measurements. Neural network-based fusion learns to combine multimodal data end-to-end, potentially capturing complex sensor interactions. Uncertainty quantification in fused estimates enables downstream systems to make risk-aware decisions.

\subsect{Edge/Embedded Deployment (Perception)} Deploying perception models on embedded hardware requires careful optimization to achieve real-time performance within resource constraints. Model compression techniques including quantization and pruning reduce computational and memory requirements. Knowledge distillation transfers capabilities from large, accurate models to smaller, efficient ones suitable for deployment. The team develops hardware-specific optimizations that exploit accelerator capabilities.

% =============================================================================
% AUTONOMY & PLANNING
% =============================================================================
\sect{Autonomy \& Planning} The Autonomy \& Planning sub-area develops the decision-making and motion planning capabilities that enable systems to act purposefully in their environments. This team translates perception outputs into actions that accomplish mission objectives while respecting safety constraints and platform limitations.

\subsect{Path Planning} Path planning computes collision-free trajectories from current position to goal locations. Global planning algorithms including A*, RRT, and PRM search for paths through environmental representations, finding routes around obstacles. Local planning refines global paths into smooth, dynamically feasible trajectories suitable for execution. Kinodynamic planning accounts for platform dynamics, generating trajectories that respect velocity, acceleration, and turning constraints. Multi-agent path planning coordinates multiple autonomous systems operating in shared environments.

\subsect{Navigation} Navigation executes planned paths in the physical world, handling discrepancies between planned and actual conditions. Waypoint following controllers track reference trajectories while compensating for disturbances and model errors. Reactive obstacle avoidance supplements planned paths with immediate responses to newly detected hazards. Terrain-aware navigation adapts behavior to surface conditions, adjusting speed and path selection based on terrain traversability. GPS-denied navigation strategies maintain operation when global positioning is unavailable.

\subsect{RL-Based Control} Reinforcement learning develops control policies through interaction with environments, potentially discovering behaviors that outperform hand-designed controllers. Policy learning for control trains neural networks to map observations directly to actions, learning from reward signals that encode desired behavior. Sim-to-real transfer addresses the challenge of deploying policies trained in simulation to physical systems. Safe reinforcement learning incorporates constraints that prevent dangerous actions during learning and deployment. Hierarchical RL architectures decompose complex tasks into manageable sub-policies.

\subsect{Mission-Level Decisions} Mission-level decision-making addresses high-level choices that span beyond immediate motion planning. Task allocation and scheduling determine what objectives to pursue and in what order, balancing priorities and resource constraints. Contingency planning prepares responses to anticipated failures or environmental changes. Multi-objective optimization balances competing goals when perfect solutions are impossible. Human-autonomy teaming integrates autonomous capabilities with human oversight.

% =============================================================================
% GENERALIZATION & ADAPTATION
% =============================================================================
\sect{Generalization \& Adaptation} The Generalization \& Adaptation sub-area ensures that systems perform reliably in novel environments they have never previously encountered. This capability is essential for autonomous systems that must operate across diverse, unpredictable conditions without manual reconfiguration.

\subsect{Domain Adaptation} Domain adaptation addresses performance degradation when models trained on source data are deployed to different target domains. Unsupervised domain adaptation techniques align feature distributions between domains without requiring target domain labels. Source-free adaptation operates without access to original training data, adapting models using only unlabeled target samples. Continual and lifelong learning enables models to incorporate new experiences without forgetting previously learned capabilities.

\subsect{Transfer Learning} Transfer learning leverages knowledge from related tasks to accelerate learning on new problems. Pre-training on large datasets followed by fine-tuning on task-specific data has proven highly effective across perception domains. Multi-task learning trains models on multiple related tasks simultaneously, learning shared representations that benefit all tasks. Cross-modal transfer applies knowledge learned from one sensor modality to another. Few-shot and zero-shot learning techniques enable recognition of new classes from minimal or no examples.

\subsect{Robustness to Novel Environments} Ensuring reliable operation in truly novel conditions requires explicit attention to robustness. Out-of-distribution detection identifies when inputs differ significantly from training data, enabling appropriate fallback behaviors. Uncertainty estimation quantifies model confidence, distinguishing reliable predictions from uncertain guesses. Graceful degradation strategies maintain safe operation as conditions exceed system capabilities, trading performance for reliability.

% =============================================================================
% SIMULATION & V&V
% =============================================================================
\sect{Simulation \& V\&V} The Simulation \& V\&V sub-area provides testing infrastructure and validation methodology that ensures system correctness before field deployment. Comprehensive testing is essential for autonomous systems where failures may have serious consequences.

\subsect{Simulation Environments} Building realistic simulation environments enables extensive testing without the cost and risk of field operations. Physics-based simulation accurately models platform dynamics, sensor responses, and environmental interactions. Sensor simulation fidelity determines how well simulation results predict real-world performance. Scenario generation creates diverse test cases systematically exploring operational conditions, edge cases, and failure modes. Digital twin development creates simulation models matched to specific physical systems.

\subsect{Testing} Systematic testing evaluates system performance across relevant conditions. Unit and integration testing verify individual components and their interactions. Hardware-in-the-loop testing connects physical hardware to simulated environments, validating hardware behavior without full field deployment. Field testing protocols structure real-world evaluation to efficiently characterize performance and identify issues. Regression testing ensures that changes do not degrade previously working capabilities.

\subsect{Validation \& Failure Analysis} Validation confirms that systems meet requirements and perform as intended. Performance metric definition establishes quantitative criteria for success, enabling objective evaluation. Edge case identification systematically discovers challenging conditions that stress system capabilities. Failure mode and effects analysis (FMEA) anticipates potential failures and their consequences. Root cause analysis investigates failures to identify underlying issues and prevent recurrence.

% =============================================================================
% INTEGRATION
% =============================================================================
\sect{Integration} The Integration sub-area combines components from all other sub-areas into complete, functioning autonomous systems. This team manages the complexity of assembling diverse subsystems into coherent wholes that operate reliably in the field.

\subsect{System Integration} Combining subsystems requires careful attention to interfaces and interactions. Software architecture design establishes the structure within which components operate, defining communication patterns and execution models. Interface definition and management ensures that components exchange information correctly, with clear specifications of data formats, timing, and error handling. Middleware selection chooses frameworks such as ROS that provide common infrastructure for component communication. Version control and CI/CD practices maintain code quality and enable rapid iteration.

\subsect{Communications \& Networking} Enabling data flow within and outside autonomous systems requires robust communication infrastructure. Internal bus architectures connect sensors, processors, and actuators with appropriate bandwidth and latency. External communication links enable remote monitoring, control, and data transmission. Bandwidth management allocates limited communication capacity among competing demands. Latency optimization reduces delays in time-critical control loops and situational awareness updates.

\subsect{Full Stack Assembly} Final system assembly brings together all components into deployable configurations. System bring-up procedures define the sequence of initialization steps required to transition from power-on to operational status. Configuration management tracks system settings and enables reproducible deployments. Deployment automation reduces manual effort and potential for human error in system preparation. Operational monitoring tracks system health during deployment, enabling early detection of issues and informed decision-making.

\end{document}
